\section{Discussion}
\label{sec:discussion}

These experiments separate four objects that are easy to collapse under the
single label \emph{token entanglement}: a behavioral association, a static
relationship between output-token rows, and a contextual computation that makes
an answer linearly readable, plus the causal influence carried by a transplanted
assistant state. They do not scale together. From Llama 8B to 70B, the static
geometry--behavior link falls decisively, the late observational trajectory
remains similar, and causal donor control moves substantially earlier. Across
tokenizers, replacing one
atomic number token with a digit sequence does not preserve the positive
association, and a seemingly reasonable normalization can reverse the sign
through sequence length.

The most direct interpretation is that the measured channel is sensitive to its
tokenizer boundary and contextually assembled inside the network. The causal
patch sharpens ``contextually assembled'': a natural donor state becomes
sufficient to make the recipient output follow donor-specific animal contrasts,
and this handoff occurs much earlier at 70B. This does not
mean that multi-token representations cannot carry information. It means that
exact sequence probability was not a drop-in replacement for the positive
atomic-token statistic in the tested models. ``Contextually assembled'' means
that a fixed output-head readout becomes progressively aligned with the saved
animal behavior at the response position. Full-state patching establishes
sufficiency, but it does not imply that a single animal vector is the causal
algorithm.

The 8B--70B dissociation constrains a simple global-geometry account. If raw
output-row similarity were a sufficient, scale-monotonic explanation, its
behavioral predictiveness and contextual causal transmission should not move in
opposite directions. Larger models can evidently preserve similar output
behavior while reorganizing which static output directions predict that behavior
and shifting donor control earlier in the forward computation. This is
compatible with activation-space, channel-location, and output-compatibility
accounts
\citep{blank2026steering,morgulis2026steering,brockers2026noise,madl2026channel},
but full-state sufficiency does not choose a unique causal feature or mechanism
among them.

The Qwen result also changes how multi-token probes should be designed. Sequence
log probability, token-normalized score, and width-controlled score answer
different statistical questions. When token count is linked to stimulus width,
per-token averaging can produce a strong and highly replicable positive result
that is entirely absent within widths. Future work should preregister a fixed
length or explicitly residualize length before interpreting a sequence-level
association as semantic or mechanistic coupling.

The main remaining bridge is training. A larger prospective study should test
whether static geometry, readout depth, or causal handoff depth predicts actual
student transfer after fine-tuning. That requires multiple teachers, student
seeds, and held-out preregistered predictions; a single cheap fine-tune would not
answer it. The present study instead supplies a tightly controlled frozen-model
mechanism result and identifies which measurements should not be substituted for
one another.
